% GENERATED FILE. Edit canonical lessons, then run npm run generate:books.
% canonical-source-hash: fnv1a64:d12082063f674ae5
% canonical-lessons: TE-C490-cedagottu, TE-C490-cirigipovu, TE-C490-masipovu, TE-C490-cilipovu, TE-C490-virigipovu

\chapter{Torn and Crushed}
\label{ch:te-torn-and-crushed}
\hlchaptermodality{490}

\begin{chapteropening}
\textbf{By the end of this chapter:} \emph{I can say to spoil, to get torn, to get grubby, to split and to break off.}

Complete the last lesson of chapter 490: I can say to spoil, to get torn, to get grubby, to split and to break off.
\end{chapteropening}

\section[\texorpdfstring{\te{చెడగొట్టు} (ceḍagoṭṭu)}{చెడగొట్టు (ceḍagoṭṭu)}]{\te{చెడగొట్టు} (ceḍagoṭṭu) — to spoil; to ruin}

\label{lesson:TE-C490-cedagottu}



\begin{quote}
Before the new one: say the Telugu for to rot, then the Telugu for to be left over.
\end{quote}

\subsection*{You'll want to know: \te{చెడగొట్టు}}
\textbf{\te{చెడగొట్టు}} — \emph{ceḍagoṭṭu} — \textquotedblleft{}to spoil; to ruin\textquotedblright{}.

\begin{quote}
\te{వర్షం} \te{మా} \te{విహారయాత్రను} \te{చెడగొట్టింది}.
\end{quote}

\begin{quote}
\emph{varṣaṁ mā vihārayātranu ceḍagoṭṭindi.}
\end{quote}

\begin{quote}
The rain spoiled our outing.
\end{quote}

Past (I …): \te{చెడగొట్టాను}, \emph{ceḍagoṭṭānu}.

\te{చెడ}, bad, + \te{కొట్టు}, to strike: to make something go bad. Said of plans, food and children who are spoilt.

\begin{cousinweb}
\textbf{\te{చెడగొట్టు}} (\emph{ceḍagoṭṭu}) is native Dravidian: \te{చెడు}, \textquotedblleft{}to go bad\textquotedblright{} (DEDR 2760), and \te{కొట్టు}, \textquotedblleft{}to strike\textquotedblright{} (DEDR 2063).
\end{cousinweb}

\begin{grammarlens}[title={What you've built}]
One more word for getting around.
\end{grammarlens}

\subsection*{Your turn}
\begin{itemize}
\raggedright
  \item \emph{Say it:} \emph{ceḍagoṭṭu}
  \item \emph{Say it:} \emph{ceḍagoṭṭu}, once more
  \item \emph{Say it:} \emph{varṣaṁ mā vihārayātranu ceḍagoṭṭindi.} — The rain spoiled our outing.
  \item \emph{Recall:} say the Telugu for to rot, then the Telugu for to be left over, then say \emph{ceḍagoṭṭu} again
\end{itemize}

\begin{tcolorbox}[breakable,colback=teal!4,colframe=teal!35!black,title={Before you move on}]
What does \te{చెడగొట్టు} mean? (\textquotedblleft{}To spoil; to ruin\textquotedblright{}.) Say it once more.
\end{tcolorbox}
\section[\texorpdfstring{\te{చిరిగిపోవు} (cirigipōvu)}{చిరిగిపోవు (cirigipōvu)}]{\te{చిరిగిపోవు} (cirigipōvu) — to get torn}

\label{lesson:TE-C490-cirigipovu}



\begin{quote}
Before the new one: say the Telugu for to be left over, then the Telugu for to spoil.
\end{quote}

\subsection*{You'll want to know: \te{చిరిగిపోవు}}
\textbf{\te{చిరిగిపోవు}} — \emph{cirigipōvu} — \textquotedblleft{}to get torn\textquotedblright{}.

\begin{quote}
\te{ముళ్ళకు} \te{తగిలి} \te{నా} \te{చొక్కా} \te{చిరిగిపోయింది}.
\end{quote}

\begin{quote}
\emph{muḷḷaku tagili nā cokkā cirigipōyindi.}
\end{quote}

\begin{quote}
My shirt caught on the thorns and got torn.
\end{quote}

Past (it …): \te{చిరిగిపోయింది}, \emph{cirigipōyindi}.

Said of cloth and paper. \te{చిరుగు} is a tear. Its subject is a thing, so the past shown is the third person: \te{చిరిగిపోయింది}, it got torn.

\begin{cousinweb}
\textbf{\te{చిరిగిపోవు}} (\emph{cirigipōvu}) is native Dravidian: DEDR 2491 lists Telugu \emph{cirugu}, \textquotedblleft{}to be torn\textquotedblright{}, beside \emph{cīru}, \textquotedblleft{}to tear\textquotedblright{}, here with \te{పోవు}, \textquotedblleft{}to go\textquotedblright{} (DEDR 4572).
\end{cousinweb}

\begin{grammarlens}[title={What you've built}]
One more word for getting around.
\end{grammarlens}

\subsection*{Your turn}
\begin{itemize}
\raggedright
  \item \emph{Say it:} \emph{cirigipōvu}
  \item \emph{Say it:} \emph{cirigipōvu}, once more
  \item \emph{Say it:} \emph{muḷḷaku tagili nā cokkā cirigipōyindi.} — My shirt caught on the thorns and got torn.
  \item \emph{Recall:} say the Telugu for to be left over, then the Telugu for to spoil, then say \emph{cirigipōvu} again
\end{itemize}

\begin{tcolorbox}[breakable,colback=teal!4,colframe=teal!35!black,title={Before you move on}]
What does \te{చిరిగిపోవు} mean? (\textquotedblleft{}To get torn\textquotedblright{}.) Say it once more.
\end{tcolorbox}
\section[\texorpdfstring{\te{మాసిపోవు} (māsipōvu)}{మాసిపోవు (māsipōvu)}]{\te{మాసిపోవు} (māsipōvu) — to get grubby; to get soiled}

\label{lesson:TE-C490-masipovu}



\begin{quote}
Before the new one: say the Telugu for to spoil, then the Telugu for to get torn.
\end{quote}

\subsection*{You'll want to know: \te{మాసిపోవు}}
\textbf{\te{మాసిపోవు}} — \emph{māsipōvu} — \textquotedblleft{}to get grubby; to get soiled\textquotedblright{}.

\begin{quote}
\te{రెండు} \te{రోజుల} \te{ప్రయాణంలో} \te{నా} \te{బట్టలు} \te{మాసిపోయాయి}.
\end{quote}

\begin{quote}
\emph{reṇḍu rōjula prayāṇaṁlō nā baṭṭalu māsipōyāyi.}
\end{quote}

\begin{quote}
After two days on the road my clothes were grubby.
\end{quote}

Past (it …): \te{మాసిపోయింది}, \emph{māsipōyindi}.

Said of clothes, walls and hands that dirt has dulled. \te{మాసిన} \te{బట్టలు}, dirty clothes. Its subject is a thing, so the past shown is the third person: \te{మాసిపోయింది}, it got grubby.

\begin{cousinweb}
\textbf{\te{మాసిపోవు}} (\emph{māsipōvu}) is native Dravidian: DEDR 4792 lists Telugu \emph{māsi}, \textquotedblleft{}dirt\textquotedblright{}, and \emph{māyu}, \textquotedblleft{}to become dirty, soiled\textquotedblright{}, here with \te{పోవు}, \textquotedblleft{}to go\textquotedblright{} (DEDR 4572).
\end{cousinweb}

\begin{grammarlens}[title={What you've built}]
One more word for getting around.
\end{grammarlens}

\subsection*{Your turn}
\begin{itemize}
\raggedright
  \item \emph{Say it:} \emph{māsipōvu}
  \item \emph{Say it:} \emph{māsipōvu}, once more
  \item \emph{Say it:} \emph{reṇḍu rōjula prayāṇaṁlō nā baṭṭalu māsipōyāyi.} — After two days on the road my clothes were grubby.
  \item \emph{Recall:} say the Telugu for to spoil, then the Telugu for to get torn, then say \emph{māsipōvu} again
\end{itemize}

\begin{tcolorbox}[breakable,colback=teal!4,colframe=teal!35!black,title={Before you move on}]
What does \te{మాసిపోవు} mean? (\textquotedblleft{}To get grubby; to get soiled\textquotedblright{}.) Say it once more.
\end{tcolorbox}
\section[\texorpdfstring{\te{చీలిపోవు} (cīlipōvu)}{చీలిపోవు (cīlipōvu)}]{\te{చీలిపోవు} (cīlipōvu) — to split; to split apart}

\label{lesson:TE-C490-cilipovu}



\begin{quote}
Before the new one: say the Telugu for to get torn, then the Telugu for to get grubby.
\end{quote}

\subsection*{You'll want to know: \te{చీలిపోవు}}
\textbf{\te{చీలిపోవు}} — \emph{cīlipōvu} — \textquotedblleft{}to split; to split apart\textquotedblright{}.

\begin{quote}
\te{భూకంపంలో} \te{గోడ} \te{రెండుగా} \te{చీలిపోయింది}.
\end{quote}

\begin{quote}
\emph{bhūkampaṁlō gōḍa reṇḍugā cīlipōyindi.}
\end{quote}

\begin{quote}
In the earthquake the wall split in two.
\end{quote}

Past (it …): \te{చీలిపోయింది}, \emph{cīlipōyindi}.

Said of wood, of walls and of a group that breaks in two. Its subject is a thing, so the past shown is the third person: \te{చీలిపోయింది}, it split.

\begin{cousinweb}
\textbf{\te{చీలిపోవు}} (\emph{cīlipōvu}) is native Dravidian: DEDR 1622 lists Telugu \emph{cīlu}, \textquotedblleft{}to break, to split, to crack\textquotedblright{}, here with \te{పోవు}, \textquotedblleft{}to go\textquotedblright{} (DEDR 4572).
\end{cousinweb}

\begin{grammarlens}[title={What you've built}]
One more word for getting around.
\end{grammarlens}

\subsection*{Your turn}
\begin{itemize}
\raggedright
  \item \emph{Say it:} \emph{cīlipōvu}
  \item \emph{Say it:} \emph{cīlipōvu}, once more
  \item \emph{Say it:} \emph{bhūkampaṁlō gōḍa reṇḍugā cīlipōyindi.} — In the earthquake the wall split in two.
  \item \emph{Recall:} say the Telugu for to get torn, then the Telugu for to get grubby, then say \emph{cīlipōvu} again
\end{itemize}

\begin{tcolorbox}[breakable,colback=teal!4,colframe=teal!35!black,title={Before you move on}]
What does \te{చీలిపోవు} mean? (\textquotedblleft{}To split; to split apart\textquotedblright{}.) Say it once more.
\end{tcolorbox}
\section[\texorpdfstring{\te{విరిగిపోవు} (virigipōvu)}{విరిగిపోవు (virigipōvu)}]{\te{విరిగిపోవు} (virigipōvu) — to break off; to snap}

\label{lesson:TE-C490-virigipovu}



\begin{quote}
Before the new one: say the Telugu for to get grubby, then the Telugu for to split.
\end{quote}

\subsection*{You'll want to know: \te{విరిగిపోవు}}
\textbf{\te{విరిగిపోవు}} — \emph{virigipōvu} — \textquotedblleft{}to break off; to snap\textquotedblright{}.

\begin{quote}
\te{గాలికి} \te{చెట్టు} \te{కొమ్మ} \te{విరిగిపోయింది}.
\end{quote}

\begin{quote}
\emph{gāliki ceṭṭu komma virigipōyindi.}
\end{quote}

\begin{quote}
A branch of the tree snapped in the wind.
\end{quote}

Past (it …): \te{విరిగిపోయింది}, \emph{virigipōyindi}.

Said of sticks, bones and chairs; glass and pots crack instead. Its subject is a thing, so the past shown is the third person: \te{విరిగిపోయింది}, it snapped.

\begin{cousinweb}
\textbf{\te{విరిగిపోవు}} (\emph{virigipōvu}) is native Dravidian: \te{విరుగు}, \textquotedblleft{}to break\textquotedblright{}, is of the family of Telugu \emph{viriyu}, \textquotedblleft{}to burst, to break open\textquotedblright{} (DEDR 5411), here with \te{పోవు}, \textquotedblleft{}to go\textquotedblright{} (DEDR 4572).
\end{cousinweb}

\begin{grammarlens}[title={What you've built}]
One more word for getting around.
\end{grammarlens}

\subsection*{Your turn}
\begin{itemize}
\raggedright
  \item \emph{Say it:} \emph{virigipōvu}
  \item \emph{Say it:} \emph{virigipōvu}, once more
  \item \emph{Say it:} \emph{gāliki ceṭṭu komma virigipōyindi.} — A branch of the tree snapped in the wind.
  \item \emph{Recall:} say the Telugu for to get grubby, then the Telugu for to split, then say \emph{virigipōvu} again
\end{itemize}

\begin{tcolorbox}[breakable,colback=teal!4,colframe=teal!35!black,title={Before you move on}]
What does \te{విరిగిపోవు} mean? (\textquotedblleft{}To break off; to snap\textquotedblright{}.) Say it once more.
\end{tcolorbox}
